% GENERATED FILE. Edit canonical lessons, then run npm run generate:books.
% canonical-source-hash: fnv1a64:36b897cb9eaf5e42
% canonical-lessons: ZH-C349-kandai, ZH-C349-xiangfa, ZH-W349-fa-rule, ZH-C349-shuofa, ZH-C349-zuofa, ZH-C349-zhuzhang

\chapter{Ways of Seeing}
\label{ch:zh-ways-of-seeing}
\hlchaptermodality{349}

\begin{chapteropening}
\textbf{By the end of this chapter:} \emph{I can ask how someone sees a question and talk about ideas and ways of doing things.}

Complete the last lesson of chapter 349: I can ask how someone sees a question and talk about ideas and ways of doing things.
\end{chapteropening}

\section[kàndài]{kàndài — to regard (in a certain way)}

\label{lesson:ZH-C349-kandai}



\begin{quote}
Before the new one: say the Mandarin for the other party, then the Mandarin for all of one mind.
\end{quote}

\subsection*{You'll want to know: kàndài}
\emph{kàndài} — \textquotedblleft{}to regard (in a certain way)\textquotedblright{}.

\begin{quote}
\emph{nǐ zěnme kàndài zhège wèntí} — How do you view this question?
\end{quote}

Its characters come later; for now it is a word for the ear.

\begin{cousinweb}
\emph{kàndài} is built from \zh{看} and \emph{dài}, both inherited from Middle Chinese, \emph{khan} and \emph{dojX} (Guangyun, in Baxter's transcription), which glosses the first 'to view' and the second 'to wait'.
\end{cousinweb}

\begin{grammarlens}[title={What you've built}]
One more word for giving reasons.
\end{grammarlens}

\subsection*{Your turn}
\begin{itemize}
\raggedright
  \item \emph{Say it:} \emph{kàndài}
  \item \emph{Say it:} \emph{kàndài}, once more
  \item \emph{Say it:} \emph{nǐ zěnme kàndài zhège wèntí} — How do you view this question?
  \item \emph{Recall:} say the Mandarin for the other party, then the Mandarin for all of one mind, then say \emph{kàndài} again
\end{itemize}

\begin{tcolorbox}[breakable,colback=teal!4,colframe=teal!35!black,title={Before you move on}]
What does kàndài mean? (\textquotedblleft{}To regard (in a certain way)\textquotedblright{}.) Say it once more.
\end{tcolorbox}
\section[\texorpdfstring{\zh{想法} (xiǎngfǎ)}{想法 (xiǎngfǎ)}]{\zh{想法} (xiǎngfǎ) — idea; way of thinking}

\label{lesson:ZH-C349-xiangfa}



\begin{quote}
Before the new one: say the Mandarin for all of one mind, then the Mandarin for to regard.
\end{quote}

\subsection*{You'll want to know: \zh{想法}}
\textbf{\zh{想法}} — \emph{xiǎngfǎ} — \textquotedblleft{}idea; way of thinking\textquotedblright{}.

\begin{quote}
\textbf{\zh{我觉得你的想法很好}} \emph{wǒ juéde nǐ de xiǎngfǎ hěn hǎo} — I think your idea is a good one.
\end{quote}

The new character, in pieces: \zh{法} — eight strokes in two pieces: on the left the water piece, three strokes, as on the left of \zh{汉}; on the right \zh{去} \emph{qù}, five strokes. The next lesson puts it in your hand.

Two third tones meet in this word: xiǎng is said with a rising tone before fǎ, and the spelling keeps both marks.

\begin{cousinweb}
\textbf{\zh{想法}} (\emph{xiǎngfǎ}) is built from \zh{想} 'to think' and \zh{法} 'a rule', both inherited from Middle Chinese, \emph{sjangX} and \emph{pjop} (Guangyun, in Baxter's transcription), which glosses the first 'to think' and the second 'a rule' and 'a norm'.
\end{cousinweb}

\begin{grammarlens}[title={What you've built}]
One more word for giving reasons.
\end{grammarlens}

\subsection*{Your turn}
\begin{itemize}
\raggedright
  \item \emph{Say it:} \emph{xiǎngfǎ}
  \item \emph{Say it:} \emph{xiǎngfǎ}, once more
  \item \emph{Say it:} \emph{wǒ juéde nǐ de xiǎngfǎ hěn hǎo} — I think your idea is a good one.
  \item \emph{Recall:} say the Mandarin for all of one mind, then the Mandarin for to regard, then say \emph{xiǎngfǎ} again
\end{itemize}

\begin{tcolorbox}[breakable,colback=teal!4,colframe=teal!35!black,title={Before you move on}]
What does \zh{想法} mean? (\textquotedblleft{}Idea; way of thinking\textquotedblright{}.) Say it once more.
\end{tcolorbox}
\section[\texorpdfstring{\zh{法} (fǎ)}{法 (fǎ)}]{\zh{法} — the water piece, then qù}

\label{lesson:ZH-W349-fa-rule}



\begin{quote}
The word first, then the character inside it.

\begin{itemize}
\raggedright
  \item \emph{Recall:} say \emph{xiǎngfǎ} — \textbf{R1}, one lesson back, before it decays
\end{itemize}

\begin{quote}
\zh{去} \zh{汉} \zh{想}
\end{quote}

\emph{Read it:} each one aloud, with nothing said first — \emph{qù}, \emph{hàn}, \emph{xiǎng}
\end{quote}

\subsection*{Script — the water piece, then qù}
\hlblockfigure{\detokenize{figures/ZH-W349-fa-rule-filmstrip.pdf}}{How \zh{法} is written, stroke by stroke}

\begin{quote}
\zh{法}
\end{quote}

\textbf{Two pieces, left before right.} Count the pieces, not the strokes.

\textbf{The left piece, the water piece, three strokes}, as on the left of \zh{汉}: a dot, a second dot under it, then a stroke that rises from the bottom left.

\textbf{The right piece, \zh{去}, five strokes}: a short horizontal, a vertical down through it, a long horizontal, then a stroke that falls left and turns right, and a dot.

3 + 5 = eight.

\subsection*{Your turn}
\begin{itemize}
\raggedright
  \item \emph{Write it:} \zh{法}, all eight, with the model in front of you
  \item \emph{Cover:} write it again from memory, and count the pieces, not the strokes
\end{itemize}

\begin{tcolorbox}[breakable,colback=teal!4,colframe=teal!35!black,title={Before you move on}]
How many pieces in \zh{法}? (\textbf{Two}.) Which character you know has the same left piece? (\textbf{\zh{汉}}, or \zh{洗}.) What is the right piece read? (\emph{qù}.)

Source: \href{https://raw.githubusercontent.com/chanind/hanzi-writer-data/68d10a4b21150cae5e1ebbd223eed289cf32d90c/data/\%E6\%B3\%95.json}{Hanzi Writer Data \zh{法}, pinned stroke paths and medians}.
\end{tcolorbox}
\section[\texorpdfstring{\zh{说法} (shuōfa)}{说法 (shuōfa)}]{\zh{说法} (shuōfa) — version of events; way of putting it}

\label{lesson:ZH-C349-shuofa}



\begin{quote}
Before the new one: say the Mandarin for to regard, then the Mandarin for a way of thinking.
\end{quote}

\subsection*{You'll want to know: \zh{说法}}
\textbf{\zh{说法}} — \emph{shuōfa} — \textquotedblleft{}version of events; way of putting it\textquotedblright{}.

\begin{quote}
\textbf{\zh{我觉得这个说法不对}} \emph{wǒ juéde zhège shuōfa bú duì} — I don't think that's the right way to put it.
\end{quote}

\begin{cousinweb}
\textbf{\zh{说法}} (\emph{shuōfa}) is built from \zh{说} 'to say' and \zh{法} 'a rule', both inherited from Middle Chinese, \emph{sywet} and \emph{pjop} (Guangyun, in Baxter's transcription), which glosses the first 'to tell' and the second 'a rule'. The second is said light here.
\end{cousinweb}

\begin{grammarlens}[title={What you've built}]
One more word for giving reasons.
\end{grammarlens}

\subsection*{Your turn}
\begin{itemize}
\raggedright
  \item \emph{Say it:} \emph{shuōfa}
  \item \emph{Say it:} \emph{shuōfa}, once more
  \item \emph{Say it:} \emph{wǒ juéde zhège shuōfa bú duì} — I don't think that's the right way to put it.
  \item \emph{Recall:} say the Mandarin for to regard, then the Mandarin for a way of thinking, then say \emph{shuōfa} again
\end{itemize}

\begin{tcolorbox}[breakable,colback=teal!4,colframe=teal!35!black,title={Before you move on}]
What does \zh{说法} mean? (\textquotedblleft{}Version of events; way of putting it\textquotedblright{}.) Say it once more.
\end{tcolorbox}
\section[\texorpdfstring{\zh{做法} (zuòfǎ)}{做法 (zuòfǎ)}]{\zh{做法} (zuòfǎ) — the way something is done; method}

\label{lesson:ZH-C349-zuofa}



\begin{quote}
Before the new one: say the Mandarin for a way of thinking, then the Mandarin for one version of events.
\end{quote}

\subsection*{You'll want to know: \zh{做法}}
\textbf{\zh{做法}} — \emph{zuòfǎ} — \textquotedblleft{}the way something is done; method\textquotedblright{}.

\begin{quote}
\textbf{\zh{我觉得这个做法不好}} \emph{wǒ juéde zhège zuòfǎ bù hǎo} — I don't think that's a good way to go about it.
\end{quote}

\begin{cousinweb}
\textbf{\zh{做法}} (\emph{zuòfǎ}) is built from \zh{做} 'to do' and \zh{法} 'a rule'. The second is inherited from Middle Chinese \emph{pjop} (Guangyun, in Baxter's transcription), glossed 'a rule'. No dictionary we use records an older form of the first, so its history is not settled.
\end{cousinweb}

\begin{grammarlens}[title={What you've built}]
One more word for giving reasons.
\end{grammarlens}

\subsection*{Your turn}
\begin{itemize}
\raggedright
  \item \emph{Say it:} \emph{zuòfǎ}
  \item \emph{Say it:} \emph{zuòfǎ}, once more
  \item \emph{Say it:} \emph{wǒ juéde zhège zuòfǎ bù hǎo} — I don't think that's a good way to go about it.
  \item \emph{Recall:} say the Mandarin for a way of thinking, then the Mandarin for one version of events, then say \emph{zuòfǎ} again
\end{itemize}

\begin{tcolorbox}[breakable,colback=teal!4,colframe=teal!35!black,title={Before you move on}]
What does \zh{做法} mean? (\textquotedblleft{}The way something is done; method\textquotedblright{}.) Say it once more.
\end{tcolorbox}
\section[zhǔzhāng]{zhǔzhāng — to argue in favour of; to advocate}

\label{lesson:ZH-C349-zhuzhang}



\begin{quote}
Before the new one: say the Mandarin for one version of events, then the Mandarin for the way something is done.
\end{quote}

\subsection*{You'll want to know: zhǔzhāng}
\emph{zhǔzhāng} — \textquotedblleft{}to argue in favour of; to advocate\textquotedblright{}.

\begin{quote}
\emph{tā zhǔzhāng dàjiā yìqǐ qù} — He's in favour of everyone going together.
\end{quote}

Its characters come later; for now it is a word for the ear.

\begin{cousinweb}
\emph{zhǔzhāng} is built from \emph{zhǔ} and \emph{zhāng}, both inherited from Middle Chinese, \emph{tsyuX} and \emph{trjang} (Guangyun, in Baxter's transcription), which glosses the first 'to be in charge' and the second 'to spread out'. How the pair came to mean arguing for something is not settled.
\end{cousinweb}

\begin{grammarlens}[title={What you've built}]
One more word for giving reasons.
\end{grammarlens}

\subsection*{Your turn}
\begin{itemize}
\raggedright
  \item \emph{Say it:} \emph{zhǔzhāng}
  \item \emph{Say it:} \emph{zhǔzhāng}, once more
  \item \emph{Say it:} \emph{tā zhǔzhāng dàjiā yìqǐ qù} — He's in favour of everyone going together.
  \item \emph{Recall:} say the Mandarin for one version of events, then the Mandarin for the way something is done, then say \emph{zhǔzhāng} again
\end{itemize}

\begin{tcolorbox}[breakable,colback=teal!4,colframe=teal!35!black,title={Before you move on}]
What does zhǔzhāng mean? (\textquotedblleft{}To argue in favour of; to advocate\textquotedblright{}.) Say it once more.
\end{tcolorbox}
